\documentclass[10pt]{article}
\usepackage[english]{babel}
\usepackage[margin=0.78in]{geometry}
\usepackage{amsmath,amssymb,booktabs,enumitem}
\usepackage[hidelinks]{hyperref}
\setlength{\parindent}{0pt}
\setlength{\parskip}{0.52\baselineskip}
\setlength{\emergencystretch}{3em}
\setlist{nosep,leftmargin=1.55em}
\urlstyle{same}
\newcommand{\USD}{\$}

\title{Post-Remediation Technical Audit of the Prospective\\
Public Emergency Codes Mortality and Economic Models}
\author{}
\date{August 8, 2026}

\begin{document}
\maketitle

\section*{Overall assessment}

Both manuscripts, their equations, tables, captions, reference lists,
machine-readable inputs, simulation programs, JSON/CSV outputs, source
transcriptions, elicitation protocol, environment record, and release validator
were re-audited after the structural revisions.

No further concrete internal mathematical, unit, sign, branch, accounting,
manuscript-to-output, or compilation defect was identified in this release.  The
remaining limitations require external work: independent expert elicitation,
linked episode data, deployment outcome data, and independent statistical
reproduction.  The manuscripts label those items as
pending and do not claim that they have been completed.

The documented command
\begin{verbatim}
python scripts/validate_release.py --rerun --compile
\end{verbatim}
passed from a temporary clean directory.  It regenerated every archived JSON
and CSV byte-for-byte, verified the SHA-256 manifest, compiled both papers in
English, and found no undefined references, undefined citations, or overfull
boxes.

\section*{High-risk technical content rechecked}

\subsection*{Estimand, unit, and notation}

\textbf{Location.} Both abstracts, executive summaries, model definitions,
result captions, quote-safe sections, conclusions, and the mature mortality
equations.

\textbf{Result.} Both papers consistently define the work as prospective and
conditional on future PEC deployment.  They state that PEC has not generated
deployment outcome data from which a PEC-specific causal effect can be
estimated and that the distributions do not represent the probability that
deployment itself will occur.

The mature master unit is a latent unique emergency episode, not a resident-year.
The factor $d_j$ is consistently described as a repeat-record and cross-pathway
allocation factor.  One person may contribute multiple episodes per year.
Pathway labels identify one final transition/reporting category rather than one
function: function lists may overlap and several functions may act within an
episode, while the episode receives one final benefit/harm/none mortality
outcome.  No conflicting initial/final, source/observer, or episode/person
descriptor was found for the frequently used symbols.

The predictive layer assigns every generated episode exactly once and asserts
conservation.  The economic model consumes the same affected-episode arrays for
the domestic mature pathways.  Near-term 988, language, and access arrays are
generated in the mortality module and passed unchanged to the economic module.

\subsection*{Clinical-effect architecture}

\textbf{Location.} Near-term 988 and CPR sections,
\path{pec_evidence_simulation.py}, mature mortality methods,
\path{pec_full_potential_simulation.py}, and
\path{pec_mortality_inputs.json}.

\textbf{Result.} The 988 external association is separated from PEC share and
causal attribution.  Both attribution quantities have explicit zero-effect
states, and the causal slab begins at zero.  The near-term gross mortality
scenario is displayed beside the mortality-neutral result and is explicitly
beneficial-only rather than net.

The CPR calculation uses timing categories and category-specific absolute
survival rates rather than a constant per-minute log-odds coefficient.
Language-to-OHCA transport includes a point mass at zero.  Building-access
mortality is zero in the primary near-term case and appears only as a
cross-endpoint sensitivity.

Positive and adverse effects have separate spike-and-slab models and separate
shared multipliers $H^+$ and $H^-$.  Gaussian and Student-$t$ copula cases,
low/medium/high dependence, four marginal forms, Tier-C-zero, access-zero,
mortality-neutral, adverse-only, and maximal-overlap cases are archived.  The
four implementation quantities are explicitly nonserial maturity indices whose
geometric mean scales an already end-to-end pathway-success probability.

The equations contain no inverse-trigonometric or coordinate-transform formula
requiring an \texttt{atan2} or branch convention.  Probability bounds,
benefit/harm exclusivity, sign conventions, and dimensionless exponents were
checked without identifying a concrete error.

\subsection*{Economic accounting}

\textbf{Location.} Economic accounting, conservative scenarios, pathway,
adoption, and sensitivity sections, and
\path{pec_economic_simulation_v2.py}.

\textbf{Result.} Public variable expenditure and capacity are separate
quantities rather than branches of a defective fixed split.  Every pathway
subtracts receiving-system burdens from public expenditure, capacity,
private-household resources, and other direct resources.  Realization
percentages apply only to positive avoided public expenditure; modeled added
public costs remain fully counted.

The primary monetary total leaves QALYs unmonetized, excludes mortality
production, and reports VSL separately.  The evidence floor removes mortality,
monetized QALYs, productivity, and property; it discounts positive direct value
but retains adverse direct-resource changes fully.

Near-term rollout scales are separate from
evidence weights.  Adoption is a stochastic rollout fraction relative to the
mature scenario.  Distribution-form, dependence, local-correlation, and Morris
results are reported.

Pathway means conserve the national resource and mortality totals.  The
feature-level linear Shapley allocations conserve value and are labeled
reporting allocations rather than causal feature effects.

\section*{Headline-output reconciliation}

The final manuscript values agree with the regenerated outputs:
\begin{itemize}
\item near-term gross mortality benefit: mean $9.4$, median $0$, P5--P95
$0$--$50.7$, with $53.9\%$ of draws exactly zero;
\item mature conditional future-deployment predictive mortality: median $214$,
P5--P95 $-67$--$1{,}319$;
\item mature direct-process evidence-floor median: \USD518 million;
\item mature mortality-neutral societal median: \USD1.921 billion; and
\item separate mature VSL-inclusive societal median: \USD4.978 billion.
\end{itemize}

These are conditional planning outputs, not observed PEC effects.  The
distribution-form results change materially, so the papers correctly avoid
presenting beta-PERT as uniquely justified.

\section*{Editorial and proofing sweep}

The bundled \path{scripts/proofing_scan.py} was run on both final PDFs.  The
mortality PDF produced no pattern hits.  The economic PDF produced two
semicolon-capitalization candidates: one is a table-cell boundary and the other
is capitalization in a bibliography title.  Both are false positives.

No verified duplicated punctuation, malformed equation-adjacent grammar,
product-name capitalization error, or inverse-tangent division branch pattern
remains.
All citation keys have matching bibliography entries.  No undefined references
or citations, malformed quotation punctuation, or verified bibliography defect
was found.  The papers contain no figures; the tables and captions agree with
the surrounding text and final machine-readable output.

\section*{External-verification status}

The following are not internal release defects, but remain prerequisites for
stronger empirical claims:
\begin{enumerate}
\item complete the deposited independent SHELF/Delphi-style elicitation protocol;
\item estimate repeat-record and cross-system episode linkage from PSAP, 988,
211, NEMSIS, hospital, and passive-monitoring data;
\item verify or replace CPR category weights reconstructed from the published
distribution and obtain OHCA-specific language-prevalence evidence;
\item collect prospective deployment outcome data before interpreting any PEC
effect causally; and
\item obtain independent statistical reproduction and written sign-off.
\end{enumerate}

\section*{Highest-priority next actions}

\begin{enumerate}
\item Freeze and publish release \texttt{pec-prospective-2026-08-08-r1}.
\item Conduct independent elicitation and linked-episode studies before
narrowing any prior or elevating a secondary modeled result.
\item Preserve the mortality-neutral, evidence-floor, and
distribution/dependence sensitivities in every subsequent revision.
\item Replace planning inputs with observed deployment measurements as they
become available; do not relabel the present distributions as observed effects.
\end{enumerate}

\end{document}